\section{Replay Buffer 管理与数据挑选}

\subsection{为何大小与采样策略关键}

Replay buffer 提供多轮训练的数据，但它的设计直接影响性能。两个核心维度是\textbf{容量}和\textbf{采样策略}：

\begin{itemize}
    \item \textbf{容量}：太小会频繁覆盖旧样本而丢失多样性，太大则降低训练时效。实务中通常选择能存下数百万条 transition。
    \item \textbf{采样策略}：均匀采样最简单，Prioritized Experience Replay 允许高 TD error 的 transition 更常被访问。
\end{itemize}

\begin{knowledgebox}{采样策略对效果的直接影响}
High TD 文本的 transition 通常包含更难的决策或崩盘场景，优先返回它们可以让 Q 函数更快收敛；但如果优先级过度，可能造成估计偏差。通常的折衷是使用 $\beta$ 衰减为 1 的 importance-sampling 权重来校正偏置。
\end{knowledgebox}

\subsection{清洗回放数据：过滤 vs. 补充}

每当 agent 执行新策略，就应评估新采集数据的质量：

\begin{itemize}
    \item 过滤：移除 reward 异常、超出边界的 transition（例如撞墙后得 reward 为 1），避免训练噪音；
    \item 补充：人为合成带有特定标签的 transition（例如 reward shaping、domain randomization）帮助 agent 学习罕见但关键的情境。
\end{itemize}

\begin{warningbox}{回放 buffer 的污染效应}
如果 buffer 长期积累低质量数据（例如 policy 崩溃后的 exit transitions），模型可能会在错误的示例上继续优化。定期清空 buffer 7\% 左右最低 reward 的样本并重新收集，可以避免“data poisoning”。
\end{warningbox}

\subsection{本章小结}

Replay buffer 的容量、采样策略和数据清洗策略共同决定 off-policy 学习的稳定性。适度的优先级采样、IS 校正以及对异常 transition 的剔除/采样，可以避免模型在旧数据上过拟合或被噪音拖累。

